\label{sec:alcance}

El estudio es una caracterización experimental del mecanismo propuesto, no una comparación general de desempeño entre arquitecturas. El objetivo no consiste en demostrar que una arquitectura supera a otra en un conjunto amplio de problemas de referencia, sino en analizar bajo qué condiciones el mecanismo de focalización logra orientar el aprendizaje hacia regiones de mayor complejidad física. Por ello, el desempeño observado en un dominio particular no se interpreta como una conclusión absoluta, sino como evidencia de comportamiento que permite delimitar condiciones de uso, límites y oportunidades de refinamiento \cite{shukla2024fair,krishnapriyan2021}.

La unidad de análisis es el comportamiento espacial de la representación aprendida. En consecuencia, la evaluación no se limita al error global del modelo, sino que considera la distribución local del error, el residual físico, la estabilidad del entrenamiento y el costo computacional asociado. El criterio detrás de esta elección es distinguir entre una mejora numérica agregada y una mejora físicamente relevante en las zonas donde la solución exige mayor capacidad representacional \cite{grossmann2024,krishnapriyan2021,shukla2024fair}.